\section{Sum Rules for Correlation Functions}
In this appendix, we prove the most general sum rule for the correlation functions at the ground state of the Kitaev model, which applies to both the pure Kitaev model and the situation with an external magnetic field (general Kitaev-like models). First, the density matrix of the ground state reads
\begin{equation}\label{eq:pure-DM}
    \rho 
    = 
    \frac{1}{2^{N-N_p}}
    \hat{\Pi}
    \sum_{[\mathfrak{D}]}
    \mathfrak{C}_{\mathfrak{D}}
    \Sigma_{\mathfrak{D}}
\end{equation}
The coefficients are defined as $\mathfrak{C}_\mathfrak{D} \equiv \eta_\mathfrak{D}\anglelr{\Sigma_\mathfrak{D}}$, where $\eta_\mathfrak{D}$ is the sign obtained by squaring the string operator $\Sigma_\mathfrak{D}^2 = \eta_\mathfrak{D}\mathds{1}$ because not every string squares to one. The projector $\hat{\Pi}$ is defined as follows,
\[
\hat{\Pi}
=
\lr{\frac{\mathds{1} + \tilde{W}_v}{2}}
\lr{\frac{\mathds{1} + \tilde{W}_h}{2}}
\prod_{p}\lr{\frac{\mathds{1} + W_p}{2}}.
\]
It projects the state to the flux-free subspace in the topological sector labeled by the expectation value of the non-contractible Wilson loops $(\tilde{W}_v, \tilde{W}_h) = (+,+)$. 

A fundamental condition for a density matrix to be a pure state is that it must square to one. Directly taking the square, we obtain 
\begin{align*}
    \rho^2 
    &= 
    \frac{1}{2^{2N-2N_p}}
    \hat{\Pi}
    \sum_{[\mathfrak{D}]}\sum_{[\mathfrak{D}']}
    \mathfrak{C}_{\mathfrak{D}}
    \mathfrak{C}_{\mathfrak{D}'}
    \varphi(\mathfrak{D},\mathfrak{D}')
    \Sigma_{\mathfrak{D}\ominus\mathfrak{D}'}\\
    &=
    \frac{1}{2^{2N-2N_p}}
    \hat{\Pi}
    \sum_{[\mathfrak{D}]}\sum_{[\mathfrak{D}']}
    \mathfrak{C}_{\mathfrak{D}}
    \mathfrak{C}_{\mathfrak{D}'}
    \frac{1}{2}
    \lrm{
    \varphi(\mathfrak{D},\mathfrak{D}')
    +
    \varphi(\mathfrak{D}',\mathfrak{D})
    }
    \Sigma_{\mathfrak{D}\ominus\mathfrak{D}'} \\
    &=
    \frac{1}{2^{2N-2N_p}}
    \hat{\Pi}
    \sum_{[\mathfrak{D}]}\sum_{[\mathfrak{D}']}
    \mathfrak{C}_{\mathfrak{D}}
    \mathfrak{C}_{\mathfrak{D}'}
    \vartheta(\mathfrak{D},\mathfrak{D}')
    \Sigma_{\mathfrak{D}\ominus\mathfrak{D}'} \\
    &= 
    \frac{1}{2^{2N-2N_p}}
    \hat{\Pi}
    \sum_{[\mathfrak{D}^c]}\sum_{[\mathfrak{D}']}
    \mathfrak{C}_{\mathfrak{D}^c\ominus\mathfrak{D}'}
    \mathfrak{C}_{\mathfrak{D}'}
    \vartheta(\mathfrak{D}^c\ominus\mathfrak{D}',\mathfrak{D}')
    \Sigma_{\mathfrak{D}^c} \\
    &=
    \frac{1}{2^{N-N_p}}
    \hat{\Pi}
    \sum_{[\mathfrak{D}]}
    \lr{
        \frac{1}{2^{N-N_p}}
        \sum_{[\mathfrak{D}']}
        \mathfrak{C}_{\mathfrak{D}\ominus\mathfrak{D}'}
        \mathfrak{C}_{\mathfrak{D}'}
        \eta_{\mathfrak{D}'}
        \vartheta(\mathfrak{D},\mathfrak{D}')
    }
    \Sigma_{\mathfrak{D}}
\end{align*}
Comparing this result with Eq.~\eqref{eq:pure-DM}, we obtain the \textit{strong condition},
\begin{equation}\label{eq:strong-condition}
    \frac{1}{2^{N-N_p}}
    \sum_{[\mathfrak{D}']}
    \eta_{\mathfrak{D}'} \vartheta(\mathfrak{D},\mathfrak{D}')
    \mathfrak{C}_{\mathfrak{D}\ominus\mathfrak{D}'}
    \mathfrak{C}_{\mathfrak{D}'}
    =
    \mathfrak{C}_{\mathfrak{D}}.
\end{equation}
The symmetric phase factor $\vartheta(\mathfrak{D},\mathfrak{D}')$ is defined by symmetrizing the phase factor $\varphi(\mathfrak{D},\mathfrak{D}')$ defined by the projective representation of the edge configuration algebra.
\[
    \vartheta(\mathfrak{D},\mathfrak{D}') 
    \equiv 
    \frac{1}{2}
    \lrm{
    1
    +
    \eta_{\mathfrak{D}}
    \eta_{\mathfrak{D}'}
    \eta_{\mathfrak{D}'\ominus\mathfrak{D}}
    }\varphi(\mathfrak{D},\mathfrak{D}')
    =
    \vartheta(\mathfrak{D}',\mathfrak{D})
    =
    \text{$0$ or $\pm 1$}
\]
Since the function $\varphi$ has the following property,
\[
\varphi(\mathfrak{D}\ominus\mathfrak{D}',\mathfrak{D}') = \eta_{\mathfrak{D}}\eta_{\mathfrak{D}\ominus\mathfrak{D}'} \varphi(\mathfrak{D}',\mathfrak{D})
=
\eta_{\mathfrak{D}'} \varphi(\mathfrak{D},\mathfrak{D}'),
\]
The symmetric phase factor obeys a similar relation
\[
\vartheta(\mathfrak{D}'\ominus\mathfrak{D},\mathfrak{D})
=
\frac{1}{2}
\lrm{
    1
    +
    \eta_{\mathfrak{D}\ominus\mathfrak{D}'}
    \eta_{\mathfrak{D}'}
    \eta_{\mathfrak{D}}
}
\varphi(\mathfrak{D}'\ominus\mathfrak{D},\mathfrak{D}) 
= 
\eta_{\mathfrak{D}}
\vartheta(\mathfrak{D}',\mathfrak{D}).
\]
Taking $\mathfrak{D}' = \varnothing$ and the property $\varphi(\mathfrak{D},\mathfrak{D}) = \eta_\mathfrak{D}$, we have
\[
\vartheta\lr{\mathfrak{D},\mathfrak{D}} = 
\frac{1}{2}\lrm{1 + \eta_\mathfrak{D}\eta_{\mathfrak{D}} \eta_{\mathfrak{D}\ominus\mathfrak{D}}} \eta_\mathfrak{D}
=
\eta_\mathfrak{D}.
\]
With this property and taking $\mathfrak{D} = \varnothing$, the strong condition reduces to the \textit{weak condition}
\begin{equation}\label{eq:weak-condition}
    \frac{1}{2^{N-N_p}}
    \sum_{[\mathfrak{D}]}
    \eta_\mathfrak{D}
    \mathfrak{C}_{\mathfrak{D}}^2
    =
    1.
\end{equation}


% \begin{lemma}
%     If $[\mathfrak{D}]\in\bo{G}$, then the following identity holds
%     \[
%     \sum_{[\mathfrak{D}']\in\bo{G}}\eta_{\mathfrak{D}\ominus\mathfrak{D}'} 
%     = 
%     \delta_{[\mathfrak{D}],[\varnothing]}
%     \]
% \end{lemma}
% \begin{proof}
%     By definition, the factor $\eta_{\mathfrak{D}}$ can be expressed as $\eta_{\mathfrak{D}} = (-1)^{\Delta(\mathfrak{D})/2}$, where $\Delta(\mathfrak{D}) \equiv n_a(\mathfrak{D})-n_b(\mathfrak{D})$. 
%     \[
%     \Delta(\mathfrak{D}\ominus\mathfrak{D}') 
%     = 
%     \Delta(\mathfrak{D})
%     +
%     \Delta(\mathfrak{D}')
%     - n_a(\mathfrak{D}\cap\mathfrak{D}') 
%     + n_b(\mathfrak{D}\cap\mathfrak{D}') 
%     \]
%     \[
%     \eta_{\mathfrak{D}\ominus\mathfrak{D}'}
%     =
%     (-1)^{\Delta(\mathfrak{D}\ominus\mathfrak{D}')}
%     =
%     \eta_{\mathfrak{D}}\eta_{\mathfrak{D}'}
%     (-1)^{n_a(\mathfrak{D}\cap\mathfrak{D}') }(-1)^{n_b(\mathfrak{D}\cap\mathfrak{D}')}
%     =
%     (-1)^{n_a(\mathfrak{D}\cap\mathfrak{D}') }(-1)^{n_b(\mathfrak{D}\cap\mathfrak{D}')}
%     %(-1)^{ \abs{\mrm{EP}(\mathfrak{D}) \cap \mrm{EP}(\mathfrak{D}')} }
%     \]
%     \[
%     \sum_{[\mathfrak{D}']\in\bo{G}} \eta_{\mathfrak{D}\ominus\mathfrak{D}'}
%     =
%     \sum_{[\mathfrak{D}']\in\bo{G}}(-1)^{n_a(\mathfrak{D}\cap\mathfrak{D}') }(-1)^{n_b(\mathfrak{D}\cap\mathfrak{D}')}
%     \]
% \end{proof}
% \begin{lemma}
%     If $[\mathfrak{D}]\in\bo{G}$, then the following identity holds
%     \[
%     \sum_{[\mathfrak{D}']\in\bo{G}} \Sigma_{\mathfrak{D}'}\Sigma_{\mathfrak{D}}\Sigma_{\mathfrak{D}'} 
%     = 
%     \delta_{[\mathfrak{D}],[\varnothing]}
%     \mathds{1}
%     \]
% \end{lemma}
% \begin{proof}
%     According to the identity $\varphi(\mathfrak{D},\mathfrak{D}') = \eta_{\mathfrak{D}}\eta_{\mathfrak{D}'}\eta_{\mathfrak{D}\ominus\mathfrak{D}'}\varphi(\mathfrak{D}',\mathfrak{D})$, we have the commutation relation 
%     \[
%     \Sigma_{\mathfrak{D}}\Sigma_{\mathfrak{D}'} = \eta_{\mathfrak{D}\ominus\mathfrak{D}'} \Sigma_{\mathfrak{D}'}\Sigma_{\mathfrak{D}}
%     \]
%     In the above identity, we have already considered $[\mathfrak{D}],[\mathfrak{D}']\in\bo{G}$ in our case. With this relation, we could rewrite the summation in the equation we wanted to prove
%     \[
%     \sum_{[\mathfrak{D}']\in\bo{G}} \Sigma_{\mathfrak{D}'}\Sigma_{\mathfrak{D}}\Sigma_{\mathfrak{D}'}
%     =
%     \sum_{[\mathfrak{D}']\in\bo{G}} 
%     \eta_{\mathfrak{D}\ominus\mathfrak{D}'}
%     \Sigma_{\mathfrak{D}}
%     \]
% \end{proof}

% Define $\ket{\Sigma_\mathfrak{D}} \equiv \ket{\mathfrak{D}}$, where $\ket{0}$ is the all-up ferromagnetic state. 
% \begin{align}\label{eq:string-overlap}
%     \braket{\mathfrak{D}'}{\mathfrak{D}}
%     &=
%     \bra{0}\Sigma_{\mathfrak{D}'}\Sigma_{\mathfrak{D}}\ket{0}
%     =
%     \varphi(\mathfrak{D}',\mathfrak{D}) 
%     \bra{0}\Sigma_{\mathfrak{D}'\ominus\mathfrak{D}}\ket{0} \nonumber \\
%     &=
%     \left\{
%     \begin{array}{lll}
%         \varphi(\mathfrak{D}',\mathfrak{D}) & ; & \text{if $\mathfrak{D}'\ominus\mathfrak{D} \subset \mathfrak{D}_{Z}$}\\
%         0 & ; & \text{otherwise.}
%     \end{array}
%     \right.
% \end{align}